\section{Related Work}
\subsection{Procedural and controllable 3D environments}
Procedural scene generators reduce the manual effort of assembling varied
training worlds, but they differ in what their controls describe. ProcTHOR
combines room layouts, assets and object placement to produce interactive houses
for embodied learning~\cite{deitke2022procthor}. Its post-generation validator
uses reachable agent positions to check access to every room. This provides a
precedent for coupling procedural generation with agent-specific reachability.
Our geometric question concerns a different setting: three-dimensional paths for a specified
spherical envelope through irregular cave boundaries represented by separate
visual and collision meshes.

Control also extends beyond selecting a random seed. Infinigen generates
natural scenes through procedural geometry and materials, exposing parameters
that determine scene content~\cite{raistrick2023infinigen}. Infinigen Indoors
adds a language for composition constraints, a solver for arranging procedural
assets, and export to real-time simulators~\cite{raistrick2024indoors}.
Holodeck translates language descriptions into scene designs and spatial
relations between retrieved objects, then optimizes their
placement~\cite{yang2024holodeck}. These approaches establish ways to connect
user intent to geometry, spatial constraints and embodied environments.
Cave Composer concentrates its controls on the cavity itself: passage structure,
local wall shape and clearance relative to a robot. Scene composition and
robot-clearance conditions are complementary forms of control.

\subsection{Underground and cave environment generation}
Cave modeling has long addressed geometry beyond smooth, uniform tunnels.
Paris et al. construct karst networks from geological conditions and
inlet--outlet information, then model conduits with implicit primitives,
blending and warping~\cite{paris2021caves}. Their method controls both network
structure and cross-section geometry. Infinigen's cave generator likewise uses
probabilistic production rules for turns, elevation changes and forks, and
combines passages with varying cross-sections through implicit
geometry~\cite{raistrick2023infinigen}. Geological coherence and geometric
variety are therefore substantive prior capabilities. Our emphasis is on
conditioning these kinds of variation on a robot's required clearance and
checking the resulting start--goal tasks; geometric resemblance to real caves
is a separate property to measure.

Robotics-oriented underground generators connect structural representations
to reusable simulation geometry. Cano et al. generate tunnel meshes from a
graph and support automated workflows across multiple simulated
environments~\cite{cano2024tunnels}. PLUME separates graph construction,
meshing and texturing, supports single- and multi-layer underground layouts,
and accepts user-supplied graphs~\cite{garcia2025plume}. Its procedural
materials, texture baking and scene export support robotic simulation across
varied passages and chambers. These systems supply important precedents for
automating underground environment creation. Cave Composer organizes its
pipeline around the relationship between that generated environment and each
sampled task: an explicit protected passage constrains generation, an occupancy
search proposes a path without consuming the construction graph, and both
mesh representations participate in acceptance. The comparison concerns how
requested morphology and clearance relate to accepted tasks on the resulting
geometry.

\subsection{Environment design and generalization}
Environment variation matters because a learner can fit its training worlds
without succeeding on held-out tasks. Procgen studies this issue with procedural
environment distributions and separate assessments of sample efficiency and
generalization~\cite{cobbe2020procgen}. Domain randomization varies rendering
conditions to support visual transfer from simulation~\cite{tobin2017randomization}.
For cave navigation, these perspectives motivate distinguishing appearance
variation from structural and geometric variation. Cave Composer assigns
appearance a separate seed, allowing texture selection to vary while a cave's
geometry and task remain fixed. Such controls define interventions for studying
environment design; a new seed alone does not define an out-of-distribution test.

Unsupervised environment design goes further by adapting the training
distribution to the learner. PAIRED uses a protagonist, an antagonist and an
environment adversary, with a regret-based objective that favors challenging,
solvable environments~\cite{dennis2020paired}. Structured, solvable task
distributions are thus an established research objective. Cave Composer instead
specifies cave geometry and task acceptance through explicit conditions and
geometric checks, independently of policy returns. It provides an environment
family within which the effects of chosen training conditions can be studied.
This places the learning question on the environment distribution: comparisons
should vary that distribution within each fixed baseline, rather than attribute
differences between learning algorithms to the generator.

\subsection{Underwater simulation}
Underwater simulators model the sensing and dynamics through which cave geometry
becomes a robot's experience. Stonefish supplies marine-robot physics, sensors
and rendering~\cite{cieslak2019stonefish}; OceanSim emphasizes GPU-accelerated,
physics-based underwater visual and acoustic sensor
simulation~\cite{song2025oceansim}. Isaac Lab provides a broader framework for
parallel robot learning~\cite{nvidia2025isaaclab}. These platforms are
complementary to Cave Composer's environment-generation role. Their vehicle and
sensor models provide the physical and perceptual conditions under which a
generated cave can be evaluated. Cave Composer supplies controlled geometry and
geometrically checked tasks for such evaluations. Studying environment design
with existing learners therefore requires a specified vehicle and observation
model in addition to geometric feasibility.
